\chapter*{Nota sull'utilizzo dei modelli di linguaggio}
\addcontentsline{toc}{chapter}{Nota sull'utilizzo dei modelli di linguaggio}

Nel corso del lavoro ho utilizzato OpenAI Codex, tramite l'estensione per Visual Studio Code, e Claude Code di Anthropic come assistenti nelle attività di sviluppo software. Li ho impiegati soprattutto per mettere in discussione le mie scelte progettuali, confrontare soluzioni alternative, individuare criticità, orientare la ricerca di documentazione tecnica da approfondire, analizzare errori e supportare la preparazione dei test. Li ho usati inoltre per controllare la coerenza tra codice, test, documentazione e artefatti sperimentali. Il loro contributo è stato particolarmente utile nei refactoring più estesi, resi necessari dalla progressiva comprensione del problema e dalla conseguente revisione dell'architettura della pipeline.

Nella scrittura della tesi ho utilizzato gli stessi strumenti come revisori critici, chiedendo loro di individuare passaggi poco chiari, incoerenze tra metodo, risultati e conclusioni, affermazioni non sufficientemente sostenute e possibili punti deboli dell'argomentazione. Li ho inoltre usati come supporto per l'organizzazione del documento, la revisione linguistica e la gestione dei sorgenti \LaTeX{}. I suggerimenti ricevuti sono stati valutati e rielaborati sulla base del codice, degli artefatti sperimentali e delle fonti consultate.

Considerata l'ampiezza delle attività affrontate tra maggio e settembre 2026, ritengo che questi strumenti abbiano avuto un ruolo importante nel rendere sostenibile il lavoro entro il tempo disponibile. Senza il loro contributo, sarebbe stato difficilmente realistico raggiungere nello stesso periodo un livello analogo di sviluppo, sperimentazione, revisione e documentazione. La scelta di intraprendere questa tesi è nata anche dal desiderio di mettermi alla prova nell'uso consapevole di tecnologie emergenti e di sperimentarne direttamente potenzialità e limiti.

L'impiego del modello multimodale locale \texttt{qwen3.5:4b} tramite Ollama per la generazione delle descrizioni delle immagini costituisce invece parte del metodo sperimentale della tesi. Configurazione, attività svolta, controlli ed esiti sono descritti nel Capitolo~\ref{cap:metodi} e discussi nel Capitolo~\ref{cap:risultati}.